% !TEX program = xelatex
\documentclass[10pt,letterpaper,twocolumn]{article}

% === GEOMETRY ===
\usepackage[top=0.75in, bottom=0.75in, left=0.75in, right=0.75in]{geometry}

% === FONTS ===
\usepackage{fontspec}
\usepackage{unicode-math}
\setmainfont{Latin Modern Roman}[
  Extension=.otf,
  UprightFont=lmroman10-regular,
  BoldFont=lmroman10-bold,
  ItalicFont=lmroman10-italic,
  BoldItalicFont=lmroman10-bolditalic,
]
\setsansfont{Latin Modern Sans}[
  Extension=.otf,
  UprightFont=lmsans10-regular,
  BoldFont=lmsans10-bold,
  ItalicFont=lmsans10-oblique,
  BoldItalicFont=lmsans10-boldoblique,
]
\setmonofont{Latin Modern Mono}[
  Extension=.otf,
  UprightFont=lmmono10-regular,
  ItalicFont=lmmono10-italic,
  Scale=0.85,
]
\newfontfamily\acbold{ywft-processing-bold}[
  Path=../../system/public/type/webfonts/,
  Extension=.ttf
]
\newfontfamily\aclight{ywft-processing-light}[
  Path=../../system/public/type/webfonts/,
  Extension=.ttf
]

% === PACKAGES ===
\usepackage{xcolor}
\usepackage{titlesec}
\usepackage{enumitem}
\usepackage{booktabs}
\usepackage{tabularx}
\usepackage{fancyhdr}
\usepackage{hyperref}
\usepackage{xurl}
\usepackage{ragged2e}
\usepackage{microtype}
\usepackage[colorspec=0.92]{draftwatermark}

% === COLORS ===
\definecolor{acpink}{RGB}{180,72,135}
\definecolor{acpurple}{RGB}{120,80,180}
\definecolor{acdark}{RGB}{64,56,74}
\definecolor{acgray}{RGB}{119,119,119}
\definecolor{acgreen}{RGB}{40,150,90}
\definecolor{draftcolor}{RGB}{180,72,135}
\definecolor{codebg}{RGB}{244,241,247}

% === DRAFT WATERMARK ===
\DraftwatermarkOptions{
  text=FIELD REPORT,
  fontsize=2.6cm,
  color=draftcolor!16,
  angle=45,
  pos={0.5\paperwidth, 0.5\paperheight}
}

% === HYPERREF ===
\hypersetup{
  colorlinks=true,
  linkcolor=acpurple,
  urlcolor=acpurple,
  citecolor=acpurple,
  pdfauthor={@jeffrey},
  pdftitle={nela computer club: Field Report on a Signal-Authenticated Raspberry Pi Commune},
}

% === SECTION FORMATTING ===
\titleformat{\section}
  {\normalfont\bfseries\normalsize\uppercase}
  {\thesection.}{0.5em}{}
\titlespacing{\section}{0pt}{1.2em}{0.3em}
\titleformat{\subsection}
  {\normalfont\bfseries\small}
  {\thesubsection}{0.5em}{}
\titlespacing{\subsection}{0pt}{0.8em}{0.2em}

% === HEADER/FOOTER ===
\pagestyle{fancy}
\fancyhf{}
\renewcommand{\headrulewidth}{0pt}
\fancyhead[C]{\footnotesize\color{draftcolor}\textit{Field report --- a friend's shell box, explored with permission}}
\fancyfoot[C]{\footnotesize\thepage}

% === LISTS / SPACING ===
\setlist[itemize]{nosep, leftmargin=1.2em, itemsep=0.1em}
\setlist[enumerate]{nosep, leftmargin=1.4em}
\setlength{\columnsep}{1.8em}
\setlength{\parindent}{1em}
\setlength{\parskip}{0.3em}
\tolerance=800
\emergencystretch=1em
\hyphenpenalty=50

\newcommand{\code}[1]{{\ttfamily\small\colorbox{codebg}{#1}}}
\newcommand{\fp}[1]{{\ttfamily\small\color{acdark}#1}}

\begin{document}

\twocolumn[{%
\begin{center}
{\acbold\fontsize{23pt}{27pt}\selectfont\color{acdark} nela computer club{\color{acpink}.}}\par
\vspace{0.35em}
{\aclight\fontsize{11pt}{13pt}\selectfont\color{acpink} Field Report on a Signal-Authenticated Raspberry Pi Commune}\par
\vspace{0.6em}
{\normalsize\href{https://prompt.ac/@jeffrey}{@jeffrey}}\par
{\small\color{acgray} Aesthetic.Computer}\par
{\small\color{acpurple} \url{https://jump.nelacomputer.club}}\par
\vspace{0.5em}
\rule{\textwidth}{1.5pt}
\vspace{0.4em}
\end{center}

\begin{center}
{\small\color{draftcolor}\textbf{[ field notes --- explored with the member's own access ]}}
\end{center}
\vspace{0.3em}

\begin{quote}
\small\noindent\textbf{Abstract.} An invitation arrived over Signal: a shell
account on the \emph{Northeast Los Angeles Computer Club}. This report documents
what the machine turns out to be --- a Raspberry~Pi~4 behind a home NAT,
published to the internet through a reverse tunnel, and gated by a bespoke SSH
gateway that authenticates members by Signal push instead of passwords. We map
its hardware, its two-hop login path, its self-hosted services, and its small
roster of inhabitants; and we record how tap-free access was wired into the
author's own machine fleet.
\end{quote}
\vspace{0.3em}
}]

\section{An invitation over Signal}

The door was a username. Point \code{ssh} at \fp{jump.nelacomputer.club},
and instead of a password prompt the banner reads:

\begin{quote}\ttfamily\small
* nela computer club *\\
Log in with Signal.\\
(...) Signal username:
\end{quote}

Type the handle, and a moment later a push notification lands on your phone:
\emph{approve this login?} Approve it, and the terminal answers
``\texttt{Signed in as "jeffrey".}'' No key exchanged in advance, no shared
secret typed --- identity is proven out-of-band, by the same app that carried
the invitation. What follows is a survey of the box on the other side of that
door, conducted with the author's own member account.

\section{The machine}

Underneath the theatrics is a \textbf{Raspberry Pi 4}. \code{uname -a} reports
the Raspberry~Pi Foundation kernel \fp{6.12.25+rpt-rpi-v8} on \fp{aarch64};
the OS is \textbf{Debian~12 (bookworm)}; the hostname is, fittingly,
\fp{computerclub}.

\begin{center}
\small
\begin{tabularx}{\columnwidth}{@{}lX@{}}
\toprule
\textbf{Field} & \textbf{Value} \\
\midrule
Board      & Raspberry Pi 4 (ARM Cortex, 4 cores) \\
Kernel     & 6.12.25+rpt-rpi-v8 (aarch64) \\
OS         & Debian GNU/Linux 12 (bookworm) \\
Memory     & $\approx$907\,MB RAM, 512\,MB swap \\
Storage    & 29\,GB SD card (\fp{/dev/mmcblk0p2}), 33\% used \\
Uptime     & short --- rebooted the morning of the visit \\
\bottomrule
\end{tabularx}
\end{center}

It is, in other words, a thirty-five-dollar single-board computer on somebody's
shelf --- and it is answering the public internet.

\section{Two hops: the gateway}

The trick to that last sentence hides in the environment. Inside the session,
\fp{SSH\_CONNECTION} reads \fp{127.0.0.1 ... 127.0.0.1 2022}. The shell you
land in did \emph{not} come from the sshd on port~22. Two custom binaries in
\fp{/usr/local/bin} explain the topology:

\begin{itemize}
\item \fp{jump-gateway} --- a $\sim$9.6\,MB service that fronts port~22, runs
the ``Log in with Signal'' exchange, stamps the member's Signal UUID into the
environment as \fp{SIGNAL\_SUBJECT}, and proxies the approved session to the
real sshd listening on \fp{127.0.0.1:2022}.
\item \fp{rathole} --- a fast reverse-proxy / NAT-traversal tunnel. The public
name \fp{jump.nelacomputer.club} resolves to a DigitalOcean address
(\fp{206.189.211.92}); \fp{rathole} stitches that relay back to the Pi living
behind a home router. The relay has the public IP; the Pi has the shell.
\end{itemize}

So a login is really two hops: \emph{you $\rightarrow$ DO relay
(rathole) $\rightarrow$ Pi gateway $\rightarrow$ internal sshd}. The Signal
approval is enforced as a PAM step in that inner sshd, which is configured
\fp{PasswordAuthentication no}, \fp{PubkeyAuthentication yes}. That last
detail is the whole reason tap-free access is possible (\S6): a valid public key
satisfies authentication before the Signal PAM prompt ever runs.

\section{What it hosts}

For a machine this small, the club runs real infrastructure:

\begin{itemize}
\item \textbf{A self-hosted git forge.} \fp{/usr/local/bin/forgejo} (a
$\sim$104\,MB Forgejo/Gitea binary, owned by \fp{max}), alongside a
\fp{/srv/gitlab} tree and a dedicated \fp{git} system user. The club keeps
its own code, not GitHub's.
\item \textbf{Physical computing.} \fp{/opt/pigpio} --- the Pi GPIO daemon.
The board is wired to hardware, not just serving bytes.
\item The usual Debian furniture otherwise; members get \fp{/bin/bash} and no
\fp{sudo}.
\end{itemize}

\section{The inhabitants}

\fp{/etc/passwd} tells a small social history. Three ``legacy'' accounts in the
1000-range predate the Signal system (no \fp{signal=} tag in their GECOS
field); a newer 3000-range cohort each carry a Signal UUID, minted by the
gateway when they joined.

\begin{center}
\small
\begin{tabularx}{\columnwidth}{@{}llX@{}}
\toprule
\textbf{User} & \textbf{uid} & \textbf{Note} \\
\midrule
max       & 1000 & founder / admin; logs in from the club LAN (\fp{10.10.10.12}); owns \fp{forgejo} \\
stalgiag  & 1001 & legacy member \\
mikey     & 1002 & legacy member \\
git       & 105  & Forgejo service account \\
max2      & 3000 & Signal-era \\
jeffrey   & 3001 & Signal-era (this author) \\
kevin     & 3002 & Signal-era \\
oioi      & 3003 & Signal-era \\
\bottomrule
\end{tabularx}
\end{center}

A commune of eight, give or take, sharing one Raspberry Pi over an encrypted
messenger. The login records show \fp{max} tending it throughout the day from
inside the house; the rest arrive from the wider internet through the tunnel.

\section{Wiring fleet access}

To reach the account from any machine in the author's fleet without a phone tap
each time, three public keys were appended --- with the member's explicit
go-ahead --- to \fp{/home/jeffrey/.ssh/authorized\_keys}: the shared
Aesthetic Computer vault key plus the per-machine ED25519 keys for the \emph{neo}
and \emph{blueberry} hosts. Because the inner sshd honours pubkey auth ahead of
the Signal PAM step, this yields an instant login:

\begin{quote}\ttfamily\small
Host club\\
\hspace*{1em}HostName jump.nelacomputer.club\\
\hspace*{1em}User jeffrey\\
\hspace*{1em}IdentityFile \textasciitilde/.ssh/id\_ed25519\\
\hspace*{1em}IdentitiesOnly yes
\end{quote}

\code{ssh club} now lands a shell on the Pi in one step from either machine. The
alias and a full operator's note live in the private vault under
\fp{nela-computer-club/}. First contact from a \emph{new} machine still uses
the Signal flow (username \fp{whistlegraph.89}, approve on phone); the keys are
strictly a convenience layer over an auth model the club owns.

\section{A note on manners}

This is somebody else's computer --- \fp{max}'s, specifically, and the club's
collectively. The exploration here used only the author's own granted account and
read-only commands; the one change made (the three convenience keys) is reversible
by deleting them from the \fp{authorized\_keys} file, and should be removed on
request if the club prefers Signal-only entry. The charm of the place is exactly
that it is small, personal, and hand-built: a Pi on a shelf in Northeast LA,
answering the whole internet, one Signal approval at a time.

\vspace{1em}
\begin{center}
{\footnotesize\color{acgray}
Explored 2026-07-07 $\cdot$ host key SHA256:VRig\ldots Y6eo\\
Aesthetic.Computer field notes}
\end{center}

\end{document}
